\loai{Loại 1: Tính nhiệt nóng chảy}
\begin{vidu}%[1P7-4.2G][Loai1]
	Biết nhiệt nóng chảy riêng của nước đá là $3,4.10^5\ J/kg$. Nhiệt lượng cần cung cấp để làm nóng chảy hoàn toàn 500 g nước đá ở $0^\circ C$ tại áp suất 1 atm là
	\choice
	{$7.10^7\ J$}
	{\True $17.10^4\ J$}
	{$17.10^5\ J$}
	{$17.10^6\ J$}
	\loigiai{
		Ta có: $Q=m.\lambda =0,5.3,4.10^5=17.10^4\ J$.
		\begin{note}
			Vì nhiệt độ $0^\circ C$ là nhiệt độ nóng chảy của nước đá rồi nên nhiệt lượng Q cấp vào sẽ chỉ dùng để làm tan nước đá mà không làm thay đổi nhiệt độ của khối nước đá.	
		\end{note}
	}
	\haidong{4}\end{vidu}

\begin{vidu}%[1P7-4.2G][Loai2]
	Một thỏi nhôm có khối lượng 1,0 kg ở $8^\circ C$.  Nhôm nóng chảy ở $658^\circ C$, nhiệt nóng chảy riêng của nhôm là $3,9.10^5\ J/kg$ và nhiệt dung riêng của nhôm là $880\ J/kg.K$. Nhiệt lượng Q cần cung cấp để làm nóng chảy hoàn toàn thỏi nhôm này là
	\choice
	{$5,9.10^6\ J$}
	{$59.10^4\ J$}
	{$4,47.10^5\ J$}
	{\True $9,62.10^5\ J$}
	\loigiai{
		Ta có: $Q=mc.\Delta t+m.\lambda \Rightarrow Q=1.880\left({658-8}\right)+1.3,9.10^5=9,62.10^5\ J$.
	}
	\haidong{6}\end{vidu}
\begin{vidu} \nguon{2}
	Phải mất bao lâu để bộ phát nhiệt $1000 \ W$ làm tan chảy $1 \ kg$ băng ở $-20^\circ C$ và áp suất 1 atm? Giả sử tất cả năng lượng từ bộ phát nhiệt đều được băng hấp thụ hoàn toàn. Nhiệt dung riêng của nước đá là $2100 \ J /kg.K$; nhiệt nóng chảy riêng của nước đá là $3,4 . 10^{5} \ J / kg$. 
	\choice
	{$4,18$ giây}
	{$41,8$ giây}
	{$5,55$ phút}
	{\True $6,37$ phút}
	\loigiai{
		\begin{itemize}
			\item Nhiệt lượng cần cung cấp là:
			\begin{align*}
				Q=m c \Delta t+\lambda m=1 . 2100 . (0-(-20))+1 . 3,4 . 10^{5}=382 . 10^3 \ J
			\end{align*}
			\item Thời gian cần thiết:
			\begin{align*}
				t=\dfrac{Q}{\mathcal{P}}=\dfrac{382 . 10^3}{1000} \approx 6,37 \ \text {phút}
			\end{align*}
		\end{itemize}
	}
\end{vidu}




\begin{vidu}%[3P6-7.2T][Loai5]
\impicinpar{Người ta dùng một laze hoạt động dưới chế độ liên tục để khoan một tấm thép. Công suất chùm laze là $\mathcal{P} = 10\ W$. Đường kính của chùm sáng là $d = 1$ mm, bề dày của tấm thép là $e = 2$ mm. Nhiệt độ ban đầu của tấm thép là $30^\circ C$. Khối lượng riêng của thép là $\rho=7800\ kg/m^3$; nhiệt dung riêng của thép là $c = 448\ J/kg.K$. Nhiệt nóng chảy riêng của thép là $\lambda = 270$ kJ/kg; nhiệt độ nóng chảy của thép là $1535^\circ C$. Coi toàn bộ năng lượng của chùm tia laze được sử dụng để làm nóng chảy lỗ khoan. Thời gian để khoan lỗ là}{\includegraphics[width=5cm]{images/8507}}
\choice
	{\True 1,157 s}
	{2,125 s}
	{2,157 s}
	{2,275 s}
	\loigiai{
			\begin{itemize}
				\item Ta có phương trình cân bằng nhiệt: $\mathcal{P}. t=mc\left( {t_{nc}}-{t_0} \right)+m. \lambda \quad (1)$.
				\item  Thể tích thép cần nung chảy hình trụ: $V=\pi \dfrac{{{d}^{2}}}{4}e$.
				\item Khối lượng của thép cần hoá lỏng: $m=\rho. V=\rho. \pi \dfrac{{{d}^{2}}}{4}e\quad (2)$
				\item Thế $(2)$ vào $(1)$ : $\mathcal{P}. t=\rho. \pi \dfrac{{{d}^{2}}}{4}ec\left({t_{nc}}-{t_0} \right)+\rho. \pi \dfrac{{{d}^{2}}}{4}e. L$.
				\item Thế số: $$\mathcal{P}. t=7800. \pi . \dfrac{{{10}^{-6}}}{4}{{. 2. 10}^{-3}}. \left[ 448. \left( 1535-30 \right)+270000 \right]=11,57
				\Rightarrow t=\dfrac{11,57}{10}=1,157\ s$$
			\end{itemize}
	}
	\haidong{10}\end{vidu}
\begin{vidu}
\impicinpar{Tính thời gian cần thiết để làm nóng chảy hoàn toàn 2 kg đồng có nhiệt độ ban đầu $30^\circ C$, trong một lò nung điện có công suất $2.10^4\ W$. Biết đồng nóng chảy ở nhiệt độ $1084^\circ C$. Biết chỉ có $50 \%$ năng lượng điện tiêu thụ của lò được dùng vào việc làm đồng nóng lên và nóng chảy hoàn toàn ở nhiệt độ không đổi. Nhiệt dung riêng của đồng là $380\ J / kg. K$; nhiệt nóng chảy riêng của đồng là $1,8.10^{5}\ J/kg$ (làm tròn kết quả đến chữ số hàng đơn vị).
}{
\begin{hinhanh}{5}
\includegraphics[width=3cm]{images/9007}
\end{hinhanh}}
	\shortans[oly]{ 116}
	\loigiai{
		
		Nhiệt lượng cần để tăng nhiệt độ là $Q=m c \Delta t=2.380.(1084-30)=801040$ J
		
		Nhiệt nóng chảy $Q_{n c}=\lambda m=1,8.10^{5}.2=3,6.10^{5}\ J$
		
		Nhiệt lượng có ích $Q_{c i}=Q+Q_{n c}=801040+3,6.10^{5}=1161040$ J
		
		Điện năng tiêu thụ $A=\dfrac{Q_{c i}}{H}=\dfrac{1161040}{0,5}=2322080\ J$
		
		Thời gian $t=\dfrac{A}{\mathcal{P}}=\dfrac{2322080}{20000} \approx 116\ s$
		
		\textbf{Trả lời ngắn:} 116
	}
	\haidong{10}\end{vidu}



\loai{Loại 2: Nóng chảy do trao đổi nhiệt}
\begin{vidu} \nguon{2}
	Trong một bình cách nhiệt, thêm $250\ g$ nước đá ở nhiệt độ $0^\circ C$ (áp suất $1\ atm$) vào $600\ g$ nước ở nhiệt độ $18^\circ C$. Biết nhiệt nóng chảy riêng của nước đá là $3,3.10^5\ J/kg$ và nhiệt dung riêng của nước là $4200\ J/kg.K$. Bỏ qua sự trao đổi nhiệt với bình chứa và môi trường.
	\begin{vdc}
		Nhiệt độ cuối cùng của hệ khi đạt trạng thái cân bằng là
		\choice
		{$18^\circ C$}
		{$2^\circ C$}
		{\True $0^\circ C$}
		{$5^\circ C$}
%		\loigiai{
%			\begin{itemize}
%				\item[A] Sai: Nếu nhiệt lượng đá tan ít, hệ mới giữ ở 18°C – không đúng
%				\item[B] Sai: Nếu toàn bộ đá tan thì nhiệt độ mới tăng lên
%				\item[C] Đúng: Lượng nhiệt từ nước không đủ làm tan hết đá → hệ dừng ở $0^\circ C$
%				\item[D] Sai: Cân bằng xảy ra sớm hơn khi hệ đạt $0^\circ C$
%			\end{itemize}
%		}

	\end{vdc}
	\dotlineEX{4}
	\begin{vdc}
		Tại thời điểm cân bằng, khối lượng nước đá còn lại là
		\choice
		{$137\ g$}
		{$100\ g$}
		{$250\ g$}
		{\True $113\ g$}
		
%		\loigiai{
%			\begin{itemize}
%				\item[A] Sai: Đây là khối lượng đã tan, không phải còn lại
%				\item[B] Sai: Tính sai lượng đá còn lại
%				\item[C] Sai: Không có trao đổi nhiệt nên đá không tan là sai
%				\item[D] Đúng: Tan được $0,137\ kg = 137\ g$ trong tổng $250\ g$ → còn lại $113\ g$
%			\end{itemize}
%		}
	\end{vdc}
	\dotlineEX{4}
\end{vidu}







\loai{Loại 3: Nóng chảy do chuyển hóa năng lượng}



\begin{vidu}
	Để xác định nhiệt nóng chảy riêng của thiếc, người ta đổ $350\ g$ thiếc nóng chảy ở nhiệt độ nóng chảy của nó là $232^\circ C$ vào $330\ g$ nước ở $7^\circ C$ đựng trong một nhiệt lượng kế có nhiệt dung bằng $100\ J / K$. Sau khi cân bằng nhiệt, nhiệt độ của nước trong nhiệt lượng kế là $32^\circ C$. Biết nhiệt dung riêng của nước là $4200\ J /kg.K$, của thiếc rắn là $230\ J /kg.K$. Bỏ qua sự trao đổi nhiệt với môi trường. Nhiệt nóng chảy riêng của thiếc là bao nhiêu J/g (làm tròn kết quả đến chữ số hàng đơn vị)?
	\shortans[oly]{60 } 
	\loigiai{
		\begin{itemize}
\item Nhiệt nóng chảy thiếc $$Q_{n c}=\lambda m=350 \lambda$$
		\item Nhiệt lượng tỏa ra để giảm nhiệt độ của thiếc $$Q_{t h}=m_{t h} c_{t h} \Delta t_{t h}=350.0,23.(232-32)=16100 J$$
		\item Nhiệt lượng thu vào của nước $$Q_{n}=m_{n} c_{n} \Delta t_{n}=330.4,2.(32-7)=34650\ J$$
		\item Nhiệt lượng thu vào của nhiệt lượng kế $$Q_{n l k}=C_{n l k} \Delta t_{n}=100.(32-7)=2500\ J$$
		\item Phương trình cân bằng nhiệt $$Q_{n c}+Q_{t h}=Q_{n}+Q_{n l k} \Rightarrow 350 \lambda+16100=34650+2500\Rightarrow \lambda \approx 60\ J /g$$
\end{itemize}
		
	}
	\haidong{9}\end{vidu}

\begin{vidu}
\impicinpar{Để đúc các vật bằng thép; người ta phải nấu chảy thép trong lò. Thép đưa vào lò có nhiệt độ $20^\circ C$. Hiệu suất của lò là $60 \%$, nghĩa là $60 \%$ nhiệt lượng cung cấp cho lò được dùng vào việc đun nóng thép. Để cung cấp nhiệt lượng, người ta đã đốt hết 200 kg than đá có năng suất tỏa nhiệt là $29.10^{6} \ J /kg$. Cho biết thép có nhiệt nóng chảy riêng là $\lambda=83,7 . 10^3 \ J /kg$; nhiệt độ nóng chảy là $1400^\circ C$; nhiệt dung riêng ở thể rắn là $0,46 \ kJ /kg$.K. Khối lượng của mẻ thép nấu chảy là bao nhiêu kg (làm tròn kết quả đến chữ số hàng đơn vị)?}{\includegraphics[width=5cm]{images/9227}} 
	
	\shortans[oly]{ 4843}
	\loigiai{
		\begin{itemize}
			\item Nhiệt lượng tỏa ra khi đốt cháy $200 \ kg$ than đá: $Q_{\text {than}}=m_{\text {than}} . q_{\text {than}}=29 . 10^{6} . 200=5,8 . 10^{9} \ J$
			\item Vì hiệu suất của lò $60 \%$ nên: $Q_{ci}=H . Q_{\text {than}}=0,6 . 5,8 . 10^{9}=3,48 . 10^{9} \ J$
			\item Nhiệt lượng cần cung cấp để thép tăng nhiệt độ lên đến $1400^\circ C$:
		\end{itemize}
		$$
		Q=m . c . \Delta t=m . 0,46 . 10^3 .(1400-20)=634800m
		$$
		\begin{itemize}
			\item Nhiệt lượng cần cung cấp để $m\ kg$ thép nóng chảy: $Q_{nc}=\lambda . m=83,7 . 10^3m$
			\item Khối lượng thép đã nấu chảy là:
		\end{itemize}
		$$
		Q_{ci}=Q+Q_{nc} \Leftrightarrow 3,48 . 10^{9}=634800 . m+83,7 . 10^3 . m \Leftrightarrow m \approx 4843\  kg.
		$$
	}
	\haidong{10}\end{vidu}
\loai{đồ thị}
\begin{vidu}
	\immini[thm]{Sự biến thiên nhiệt độ của khối nước đá đựng trong ca nhôm theo nhiệt lượng cung cấp được cho trên đồ thị. Cho nhiệt dung riêng của nước và nhôm lần lượt là $c_1=4200 \ J / kg . K$; $c_2=880 \ J / kg . K$, nhiệt nóng chảy riêng của nước đá là $\lambda=3,4 . 10^{5} \ J / kg . K$. Ca nhôm có khối lượng bao nhiêu kg (làm tròn kết quả đến chữ số hàng phần trăm)?
}{
\begin{hinhanh}{7}
		\begin{tikzpicture}[scale=1,draw=Mapcolor, very thick, >=stealth']
			\tkzDefPoints{0/0/o, 5.5/0/x, 0/3/y, 0/2/a, 4.5/2/b, 4.5/0/c, 4.25/0/d}
			\tkzDrawSegments[very thick, Mapcolor,->](o,x o,y)
			\tkzLabelPoint[below left](o){$O$}
			\tkzLabelPoint[above](x){$Q\ (kJ)$}
			\tkzLabelPoint[above](y){$t^\circ C$}
			\tkzDrawSegments[dashed,thin](a,b b,c)
			\tkzDrawSegments[very thick,blue](o,d b,d)
			\tkzDrawPoints[size=2,fill=black](o,a,c)
			\tkzDrawPoints[size=3,fill=blue](d)
			\tkzLabelPoint[left](a){$2$}
			\tkzLabelPoint[below left](d){$170$}
			\tkzLabelPoint[below right](c){$175$}
	\end{tikzpicture}\end{hinhanh}}	\shortans[oly]{ 0,45}
	\loigiai{
		\begin{itemize}
			\item \textbf{Trả lời ngắn: 0,45}
			\item Dựa vào đồ thị, nhiệt lượng nước đá thu vào để nóng chảy hoàn toàn ở $0^\circ C$ là $170 \ kJ$. Trong quá trình này, ca nhôm không nhận nhiệt.
			\item Khối lượng nước đá là: $Q_1=m_1 . \lambda \Leftrightarrow m_1=\dfrac{Q_1}{\lambda}=\dfrac{170000}{3,4 . 10^{5}}=0,5\ kg$.
			\item Nhiệt lượng nước đá và ca nhôm thu vào để tăng nhiệt độ từ $0^\circ C$ đến $2^\circ C$ là:
		\end{itemize}
		$$
		175\ kJ-170\ kJ=\left(m_1 . c_1+m_2 . c_2\right) .(2-0) \Leftrightarrow 5000=\left(0,5 . 4200+m_2 . 880\right) . 2 \Leftrightarrow m_2 \approx \dfrac{5}{11}\ kg.
		$$
	}
	\haidong{12}\end{vidu}
	\begin{vidu}
		\immini[thm]{Trong một bình cách nhiệt có: $1\ kg$ nước đá, $1\ kg$ chất $A$ dễ nóng chảy, không tan trong nước và một bếp điện có công suất không đổi, nhiệt dung riêng không đáng kể. Nhiệt độ ban đầu trong bình là $-40^\circ C$. Sau khi cho bếp hoạt động, nhiệt độ trong bình biến đổi theo thời gian như đồ thị trong hình bên. Cho biết nhiệt dung riêng của nước đá là $c_{\text{đ}} = 2.10^3\ J/kg.K$, của chất rắn $A$ là $c = 10^3\ J/kg.K$.}{
\begin{hinhanh}{7}\begin{tikzpicture}[scale=1,draw=Mapcolor, very thick, >=stealth']
				\tkzDefPoints{0/0/o, 5.5/0/x, 0/1.5/y, 0/-2/y', 1/0/a, 2.67/0/b, 4/0/p, 0/-1/d, 1/-1/m, 2.67/-1/n}
				\tkzDrawSegments[very thick,Mapcolor,->](o,x y',y)
				\tkzLabelPoint[left](o){$O$}
				\tkzLabelPoint[above](x){$\tau$ (phút)}
				\tkzLabelPoint[above](y){$t^\circ C$}
				\tkzDrawSegments[dashed](a,m b,n d,m)
				\tkzDrawSegments[very thick, blue](y',m m,n n,p)
				\tkzDrawPoints[size=3,fill=white](o,y',a,b,p,d,m,n)
				\tkzLabelPoint[left](d){$-20$}
				\tkzLabelPoint[left](y'){$-40$}
				\tkzLabelPoint[below](m){$M$}
				\tkzLabelPoint[below](n){$N$}
				\tkzLabelPoint[above](a){$1$}
				\tkzLabelPoint[above](b){$\dfrac{8}{3}$}
				\tkzLabelPoint[above](p){$4$}
				\tkzLabelPoint[below](p){$P$}
				\tkzLabelPoint[below right](y'){$A$}
			\end{tikzpicture}
\end{hinhanh}}
		\begin{vdc}
			Nhiệt nóng chảy riêng của chất rắn A là
			\choice
			{$10^3\ J/kg$}
			{$6.10^4\ J/kg$}
			{\True $10^5\ J/kg$}
			{$1,1.10^5\ J/kg$}
%			\loigiai{
%				\begin{itemize}
%					\item[A] Sai: Đây là nhiệt dung riêng, không phải nhiệt nóng chảy
%					\item[B] Sai: Tính sai năng lượng trong khoảng thời gian nóng chảy
%					\item[C] Đúng: $\lambda_A = \dfrac{Q_{MN}}{m_A} = \dfrac{\mathcal{P} \cdot \tau_{MN}}{1} = 1000 \cdot 100 = 10^5\ J/kg$
%					\item[D] Sai: Giá trị cao hơn thực tế tính được
%				\end{itemize}
%			}
\dotlineEX{4}
		\end{vdc}
		
		\begin{vdc}
			Nhiệt dung riêng của chất A sau khi đã chảy lỏng là
			\choice
			{$1000\ J/kg.K$}
			{$1500\ J/kg.K$}
			{\True $2000\ J/kg.K$}
			{$2500\ J/kg.K$}
%			\loigiai{
%				\begin{itemize}
%					\item[A] Sai: Là nhiệt dung riêng của chất A khi còn rắn
%					\item[B] Sai: Tính thiếu nhiệt lượng hấp thụ trong pha lỏng
%					\item[C] Đúng: Áp dụng phương trình $Q = (m_1 c_1 + m_2 c_2) \Delta t$, thay số và giải được $c = 2000\ J/kg.K$
%					\item[D] Sai: Lớn hơn giá trị thực tế tính được
%				\end{itemize}
%			}
\dotlineEX{4}
		\end{vdc}
	\end{vidu}
\loai{Loại 5: Bài toán nâng cao}
\begin{vidu} \nguon{2}
	Trong một bình thành mỏng thẳng đứng diện tích đáy $S=100 \ cm^2$ chứa nước và nước đá ở nhiệt độ $t_1=0^\circ C$, khối lượng nước gấp 10 lần khối lượng nước đá. Một thiết bị bằng thép được đốt nóng tới $t_2=$ $80^\circ C$ rồi nhúng ngập trong nước, ngay sau đó mức nước trong bình dâng lên cao thêm $h=3 \ cm$. Biết rằng khi trạng thái cân bằng nhiệt được thiết lập trong bình nhiệt độ của nó là $t=5^\circ C$. Bỏ qua sự trao đổi nhiệt với bình và môi trường. Cho biết nhiệt dung riêng của nước là $4200 \ J / kg . K$, của thép là $500 \ J /kg$.K. Nhiệt nóng chảy riêng của nước đá là $330 \ kJ /kg$, khối lượng riêng của thép là $7700 \ kg / m^3$. Khối lượng của nước lúc đầu trong bình bằng bao nhiêu kg (làm tròn kết quả đến chữ số hàng phần trăm)?
	\shortans[oly]{1,54 } 
	\loigiai{
		\begin{itemize}
			\item Trả lời ngắn: 1,54
			\item Gọi khối lượng nước đá trong bình lúc đầu là $m_0$ thì khối lượng nước trong bình là $10 \ m_0$
			\item Thể tích của khối thép đúng bằng thể tích nước dâng lên:
			\begin{align*}
				V_{t}=h. S=3.100=300 \ cm^3=0,3.10^{-3} \ m^3
			\end{align*}
			\item Khối lượng của khối thép: 
			\begin{align*}
				\quad m_{t}=\rho_{t} . V_{t}=0,3 . 10^{-3} . 7700=2,31 \ kg
			\end{align*}
			\item Phương trình cân bằng nhiệt:
			\begin{align*}
				m_{t} c_{t}\left(t_2-t\right)=m_0 \lambda+\left(m_0+10 m_0\right) c_{n}\left(t-t_1\right) \Rightarrow m_0=0,154 \ kg \Rightarrow m=10. m_0=1,54 \ kg
			\end{align*}
		\end{itemize}
	}
\end{vidu}
\begin{vidu} \nguon{2}
	Một chiếc cốc hình trụ, khối lượng $m$, trong đó chứa một lượng nước cũng có khối lượng bằng $m$ đang ở nhiệt độ $t_1=10^\circ C$. Người ta thả vào cốc một lượng nước đá khối lượng $M$ đang ở nhiệt độ $0^\circ C$ thì cục nước đá đó chỉ tan được $\dfrac{1}{3}$ khối lượng của nó và luôn nổi khi tan. Rót thêm một lượng nước có nhiệt độ $40^\circ C$ và có khối lượng $(m+M)$. Khi cân bằng nhiệt thì nhiệt độ của cốc lại là $10^\circ C$. Bỏ qua sự trao đổi nhiệt với môi trường xung quanh, sự giãn nở vì nhiệt của nước và cốc. Biết nhiệt dung riêng của nước là $c=4200 \ J /kg.K$, nhiệt nóng chảy riêng của nước đá là $\lambda=336000 \ J /kg$. Nhiệt dung riêng của chất làm cốc là 
	\choice
	{$1300 \ J / kg . K$}
	{\True $1400 \ J / kg . K$}
	{$1380 \ J / kg . K$}
	{$1210 \ J / kg . K$}
	
	\loigiai{
		\begin{itemize}
			\item Do nước đá không tan hết nên khi cân bằng nhiệt thì hệ có nhiệt độ $0^{\circ} \mathrm{C}$
		\item Phương trình cân bằng nhiệt diễn tả quá trình cục nước đá tan một phần ba là:
		\begin{align*}
		\dfrac{M}{3} . \lambda=\mathrm{m}\left(\mathrm{c}+\mathrm{c}_1\right) .(10-0)\Rightarrow\frac{M}{3} . \lambda=\mathrm{m}\left(\mathrm{c}+\mathrm{c}_1\right) . 10
	\end{align*}		
		\item Ta có phương trình cân bằng nhiệt:
		\begin{align*}
			& \dfrac{2 M}{3} . \lambda+M c(10-0)+m\left(c+c_1\right) .(10-0)=(M+m) . c .(40-10) \\
			\Leftrightarrow & \dfrac{2 M}{3} . \lambda+10 M c+10 m\left(c+c_1\right)=30(M+m) c \\
			\Leftrightarrow & \left(\dfrac{2}{3} . \lambda-20 c\right) M=m\left(2 c-c_1\right) 10
		\end{align*}
		\item Từ (1) và (2) ta có: \begin{align*}
		&\dfrac{\lambda}{2 \lambda-60 c}=\dfrac{c+c_1}{2 c-c_1}\Rightarrow 60 c^2=(3 \lambda-60 \mathrm{c}) . \mathrm{c}_1\\
		\Rightarrow &c_1=\dfrac{20 c^2}{\lambda-20 c}=1400 \mathrm{~J} / \mathrm{kg} \text {.độ }
	\end{align*}
		\end{itemize}
	}
\end{vidu}